\documentclass{article}

\setlength{\headsep}{0.75 in}
\setlength{\parindent}{0 in}
\setlength{\parskip}{0.1 in}

%=====================================================
% Add PACKAGES Here (You typically would not need to):
%=====================================================

\usepackage{xcolor}
\usepackage[margin=1in]{geometry} % Used for page margins, often replaces fullpage
\usepackage{amsmath,amsthm,amssymb} % amssymb added for common math symbols

\usepackage{enumitem} % For custom list environments

\usepackage{tcolorbox}
% \usepackage{fullpage} % Geometry handles this better
\usepackage[parfill]{parskip} % For paragraph spacing
\usepackage{hyperref}
\usepackage[nameinlink, capitalise]{cleveref} % For smart referencing

\usepackage{algorithm2e} % For algorithms
\usepackage{tikz} % For drawing diagrams
\usepackage{siunitx} % For units
\usepackage{graphicx}
\usepackage{booktabs}

\usetikzlibrary{arrows.meta, positioning}

\crefname{algocf}{Algorithm}{Algorithms}

\RestyleAlgo{ruled} % Style for algorithms

\definecolor{linkc}{rgb}{0.1, 0.5, 0.7}
\definecolor{citec}{rgb}{0.6, 0.3, 0.7}
\definecolor{urlc}{rgb}{0.5, 0.1, 0.2}
\hypersetup{
    colorlinks=true,
    linkcolor=linkc,
    citecolor=citec,
    urlcolor=urlc
}

\renewcommand{\d}{\,\mathrm{d}} % Differential d


%=====================================================
% Ignore This Part (But Do NOT Delete It:)
%=====================================================

\theoremstyle{definition}
\newtheorem{problem}{Problem}
\newtheorem*{fun}{Fun with Algorithms}
\newtheorem*{challenge}{Challenge Yourself}
\def\fline{\rule{0.75\linewidth}{0.5pt}}
\newcommand{\finishline}{\begin{center}\fline\end{center}}
\newtheorem*{solution*}{Solution}
\newenvironment{solution}{\begin{solution*}}{{\finishline} \end{solution*}}
\newcommand{\grade}[1]{\hfill{\textbf{($\mathbf{#1}$ points)}}}
\newcommand{\thisdate}{\today}
\newcommand{\thissemester}{\textbf{Rutgers: Spring 2026}}
\newcommand{\thiscourse}{CS 344: Algorithms} 
\newcommand{\thishomework}{Number} 
\newcommand{\thisname}{Name} 
\newcommand{\thisextension}{Yes/No} 

\headheight 40pt              
\headsep 10pt


\newcommand{\thisheading}{
   \noindent
   \begin{center}
   \framebox{
      \vbox{\vspace{2mm}
    \hbox to 6.28in { \textbf{\thiscourse \hfill \thissemester} }
       \vspace{4mm}
       \hbox to 6.28in { {\Large \hfill Homework \#\thishomework \hfill} }
       \vspace{2mm}
         \hbox to 6.28in { { \hfill \hfill} }
       \vspace{2mm}
       \hbox to 6.28in { \emph{Name(s): \thisname \hfill }}
      \vspace{2mm}}
      }
   \end{center}
   \bigskip
}

%=====================================================
% Some useful MACROS (you can define your own in the same exact way also)
%=====================================================


\newcommand{\ceil}[1]{{\left\lceil{#1}\right\rceil}}
\newcommand{\floor}[1]{{\left\lfloor{#1}\right\rfloor}}
\newcommand{\prob}[1]{\Pr\paren{#1}}
\newcommand{\expect}[1]{\Exp\bracket{#1}}
\newcommand{\var}[1]{\textnormal{Var}\bracket{#1}}
\newcommand{\set}[1]{\ensuremath{\left\{ #1 \right\}}}
\newcommand{\poly}{\mbox{\rm poly}}


%=====================================================
% Fill Out This Part With Your Own Information:
%=====================================================


\renewcommand{\thishomework}{5} %Homework number
\renewcommand{\thisname}{Ashlyn$_1$ Balicki$_1$, Iris$_2$ Martinez$_2$, May$_3$ Stepper$_3$} % Enter your name here


\begin{document}

\thisheading

\subsection*{Homework Policy}
\begin{itemize}
\item You may consult all course materials (lecture notes, textbook, slides, etc.) while writing your solutions. No external resources are permitted.


\item \textbf{All submissions should be typeset in LaTeX, handwritten solutions are not accepted}. 

\item Remember to always \textbf{prove the correctness} of your algorithms and \textbf{analyze their running time} (or any other efficiency measure asked in the question).

\item Homework assignments are completed and submitted in fixed groups of 3 students. Only \textbf{one member of the group submits} the full solution on Canvas. The submission must include the names of all group members. The other two members do \textbf{not need to submit anything separately}.

\item You may discuss the problems with any of your classmates, even those outside your group. However, \textbf{each group must write and submit their solutions independently}.

\item Attempt the challenge problems only if you find them interesting. 
\end{itemize}



\subsection*{Acknowledgment:} 
\begin{tcolorbox}[colframe=black,colback=white,arc=5mm]
Our team members, Ashlyn Balicki (zab30), Iris Martinez (bam322), and May Stepper (), collaborated extensively on this work. \textbf{We hereby affirm that all external assistance utilized in the preparation of this document has been properly acknowledged}.
\end{tcolorbox}

\finishline
\newpage


\newpage







\begin{problem}
A graph \( G = (V, E) \) is dense if \( |E| = \Theta(V^2) \). Using binary heaps, Prim's algorithm for finding a minimum spanning tree runs in \( O(E \log V) \). With Fibonacci heaps, the runtime improves to \( O(E + V \log V) \), which achieves linear time $O(E)$ for dense graphs.

Describe an alternative algorithm that:
\begin{itemize}
    \item Uses a simple data structure (avoiding Fibonacci heaps).
    \item Runs in linear time $O(E)$for dense graphs.
  \end{itemize}

\end{problem}


\begin{solution}
    \

    \textbf{Algorithm: }
    \
     
    Pick an arbitrary vertex in the graph. Now, enqueue the edges linked to that vertex. Extract-Min, giving you the cheapest edge from your current vertex.

    Mark the vertex as visited.

    Continue until there are no more unvisited nodes.

    For the priority queue, use an array. Enqueue by appending new edges to the end of the array. Extract-min by performing a linear search for the smallest item in the array.
    \

    \textbf{Proof: }
    \

    Let T be the set of edges in the MST generated with this algorithm. Let X be an alternative set of edges that is different from T that produces a cheaper MST. 

    In this case, if $X_i$ is the first different vertex and $X_i > T_i$ such that $X_{i+1} < T_{i+1}$ and $X_i + X_{i+1} < T_i + T_{i+1}$, then our algorithm would simply have dequeued that cheaper edge directly, rather than going through the more expensive edge. Therefore our algorithm always finds the MST. 
    \

    \textbf{Runtime: }
    \
    
    Using an array for a priority queue takes O(V) time per extract-min. We run that a maximum of V times, once for each vertex. Therefore, our algorithm runs in $O(V^2)$ time, which is the same as $O(E)$ because $|E| = \Theta(V^2)$

\end{solution}

\newpage


\begin{problem}
Demonstrate failure of the following algorithm on a directed graph \( G \) with edge weights in \( \{-1, 0, 1\} \). The algorithm works by transforming \( G \) into a new graph \( G' \), where all edge weights are increased by 1, i.e., for every edge \( (x, y) \) in \( G \), the new weight in \( G' \) is \( w'(x, y) = w(x, y) + 1 \). The algorithm then runs Dijkstra’s algorithm on \( G' \) starting at source \( s \) to compute the shortest \( s \)-\( t \) path. Finally, the algorithm outputs the weight of this path in the original graph \( G \).

To show that this algorithm fails, provide the following details:

\begin{itemize}
    \item The graph \( G \), including its vertices, edges, and the original edge weights \( w(x, y) \). Ensure \( G \) contains at most 7 vertices.
    \item The source vertex \( s \) and the target vertex \( t \).
    \item The transformation to \( G' \), including all edge weights \( w'(x, y) = w(x, y) + 1 \).
    \item The priority queue values during the execution of Dijkstra’s algorithm on \( G' \), showing the shortest path computation step by step.
    \item The actual shortest distance from \( s \) to \( t \) in the original graph \( G \).
    \item The incorrect shortest distance from \( s \) to \( t \) returned by the algorithm, demonstrating where and why it fails.
\end{itemize}



\end{problem}

\begin{solution}

\

Graph $G$:
\begin{align*}
s&: a (-1), t(0) \\
a&: b(-1) \\
b&: t(1)
\end{align*}

Graph $G'$:
\begin{align*}
s'&: a'(0), t'(1) \\
a'&: b'(0) \\
b'&: t(2)
\end{align*}

Dijkstra's Algorithm on $G'$: \\
$s'a', a'b', s't', b't'$

\begin{center}
\begin{tabular}{ll}
\toprule
Vertex & Weight (Previous) \\
\midrule
$s'$    & $0$               \\
$a'$    & $0 (s')$          \\
$b'$    & $0 (a')$          \\
$t'$    & $1 (s')$          \\
\bottomrule
\end{tabular}
\end{center}

The algorithm then checks $bt$, which gives a worse path ($2$) than the existing path.

The actual shortest distance from $s$ to $t$ in $G$ is $s,a,b,t$, with a weight of $-1$. The algorithm returns $st$, which has a weight of $0$ in $G$. This fails, because by adding $1$ to every edge, the algorithm is adding $l$ to each path, where $l$ is the length of the path. This causes a bias for shorter paths, which might not be the fastest.
\end{solution}

\newpage



\begin{problem}
You are given a weighted undirected graph $G=(V, E)$ with integer weights $w_e \in\{1,2, \ldots, W\}$ on each edge $e$, where $W \leq 25$. Given two vertices $s, t \in V$, the goal is to find the minimum weight path (or shortest path) from $s$ to $t$. Recall that Dijkstra's algorithm solves this problem in $O(n+m \log m)$ time even if we do not have the condition that $W \leq 25$. However, we now want to use this extra condition to design an even faster algorithm.

Design and analyze an algorithm to find the minimum weight (shortest) $s-t$ path on these restricted weighted graphs in $O(V+E)$ time.



\end{problem}


\begin{solution}

    \

    \textbf{Algorithm:}

    \

    Initialize a boolean array for tracking visited vertices.
    
    Initialize an array of length $25V$, buckets.

    Initialize an array of distances of size V. Initialize distances[s] to 0 rather than $\infty$

    Put the starting vertex s into a linked list node at position 0 of the buckets array.

    Set a pointer variable curDist to 0.

    Set a value num elements = 1.

    Initialize a table that matches nodes to their ancestors to reconstruct the path.
    
    While buckets isn't empty(num elements $> 0$), find the next nonempty bucket in the buckets array( linear search the array with the variable curDist for an index with a node in it).

    Pick the first unvisted node in that bucket to explore next. 

    Calculate the tentative distance to each unvisited node accessible from the current node.

    Update those distances in the distance array if they are less than the current values. Move the nodes into the corresponding places in the buckets array based on the new distance array values

    If we update any distances, update the prev table to match the current vertex to the previous as well.

    When the priority queue is empty, follow the sequence in the prev table to find the minimum weight path from $s$ to any other vertex.
    

    \textbf{Proof: }
    \

    Fundamentally this is the same as dijkstra's. We can prove that our approach with an array of linked lists is functionally the same as a priority queue. By incrementing curDist from 0 up to the maximum possible distance of $25V$, we always extract-min. The bucket placement of a given node is the same as the priority, and shuffling nodes across the buckets array is the same as updating priority in a priority queue. 

    Because this approach is the same as having a normal priority queue, dijkstra's doesn't change and we have a correct algorithm.
    \

    \textbf{Runtime: }
    \

    The runtime analysis for this is the same as standard dijkstra's, however: our usage of an array of buckets allows us to update priorities in constant time. Since we only visit every vertex once, we only work on every edge once too. This results in E cost updates taking O(E) time.  
    Extract min is still constant time(in this case this is the distance-based sweep of the buckets array), which means that it also still takes O(V) time to find each next vertex to check.
    
    Therefore, dijkstra's modified in this way runs in $O(V + E)$ when $W\leq25$
\end{solution}

\newpage





\begin{problem}
    We consider two types of graphs:
\begin{itemize}
    \item \textbf{Vertex-Weighted Graph} \( H = (V, E) \): Each vertex \( v \) has a weight \( c(v) \). The weight of a path is the sum of vertex weights along the path. Let \( \text{dist}_H(s, v) \) represent the shortest path from \( s \) to \( v \).
    \item \textbf{Edge-Weighted Graph} \( G = (V, E) \): Each edge \( (u, v) \) has a weight \( w(u, v) \). Let \( \text{dist}_G(s, v) \) represent the shortest path from \( s \) to \( v \).
\end{itemize}

\begin{enumerate}
    \item[a)]


Suppose you have an algorithm \( \text{EdgeWeightedSP}(G, s) \) that solves the edge-weighted shortest path problem in \( O(|E(G)|) \). Write pseudocode for \( \text{VertexWeightedSP}(H, s) \), which solves the vertex-weighted shortest path problem in \( O(|E(H)|) \). You can use \(\text{EdgeWeightedSP}(G, s) \) as a blackbox.


\begin{solution}
    \

    \textbf{Algorithm: }
    \
    
    Run a modified BFS on on the graph H starting from an arbitrary vertex v. Rather than marking vertices as visited, we mark edges as visited. This makes sure we hit every possible edge.

    Every time we follow an edge $E(u,v)$, we create an edge in a graph $H'$ with the weight of the destination vertex $c(v)$

    Now we can run EdgeWeightedSP() to solve the verted weighted problem 

    Lastly, we add the weight of the origin vertex to our calculated distance for our final answer.
    \

    \textbf{Proof:}
    Given that we are using EdgeWeighted as a black box, we only need to prove that our transformation of the vertex weighted graph into an edge-weighted one is accurate.
    \

    By running BFS, we make sure to follow every edge.

    By making the edge weighted graph edge weight equal to the weight of the destination vertex, we are encoding how much that vertex adds to the weight. This is the same as if it were a vertex weighted graph. 

    Because this is the same, now the edgeweightedSP function can solve the vertex weight problem. All we need to do at the end is add the weight of the origin vertex, which gets lost in the conversion to an edge weighted graph

    \textbf{Runtime:}
    The creation of our edge graph from our vertex graph takes $O(E)$ time. The blackbox alg then works in $O(E(G)) = O(E(H))$, therefore our algorithm runs in $O(E(H))$
\end{solution}

\newpage




\item[b)] 
Suppose you have an algorithm \( \text{VertexWeightedSP}(H, s) \) that solves the vertex-weighted shortest path problem in \( O(|E(H)|) \). Write pseudocode for \( \text{EdgeWeightedSP}(G, s) \), which solves the edge-weighted shortest path problem in \( O(|E(G)|) \). You can use \( \text{VertexWeightedSP}(H, s) \) as a blackbox.

\end{enumerate}

\end{problem}


\begin{solution}

    \
    \textbf{Algorithm: }
    \
    
    We're starting from an edge weighted graph G.
    
    We run BFS modified to visit all edges as we did in part (a).

    This time, however, every time we visit a new edge, we create an intermediate vertex of the same weight of the edge we visited and connect it to two new copies of the origin and destination vertices.

    When we're done, we have a graph G'.

    We can now run the VertexWeightedSP algorithm on G' to solve the edgeweightedSP problem.

    \

    \textbf{Proof:}
    Given that we are using VertexWeightedSP as a black box, we only need to prove that our transformation of the edge weighted graph into an vertex-weighted one is accurate.
    \

    By running BFS, we make sure to follow every edge.

    We encode the weight of each edge by creating a vertex with that weight along the path that edge traces in our new graph.

    Because the weight from u to v is the same now, we can use VertexWeightedSP to solve the shortest path problem successfully.
    \textbf{Runtime:}
    The creation of our vertex-weighted graph takes $O(E)$ time. The blackbox alg then works in $O(E(H)) = O(E(G))$, therefore our algorithm runs in $O(E(G))$
\end{solution}

\newpage




\begin{problem}
    Given a directed graph \( G = (V, E) \) with positive integer edge weights, let \( s \) be the source vertex. Assume all shortest path distances satisfy \( \text{dist}(s, v) \leq B \) for some parameter \( B \). Design a data structure \( D \) such that running Dijkstra’s algorithm using \( D \) results in a runtime of \( O(|E| + B) \).


\begin{itemize}
    \item Describe the data structure \( D \) and its three operations:
        \begin{itemize}
            \item \textbf{Insert:} Add a vertex with its priority.
            \item \textbf{Extract-Min:} Remove and return the vertex with the smallest priority.
            \item \textbf{Decrease-Key:} Update the priority of a vertex.
        \end{itemize}
    \item Write pseudocode for these operations to show how \( D \) achieves \( O(1) \) or \( O(B) \) time per operation as appropriate.
    \item Justify why the runtime of Dijkstra’s algorithm becomes \( O(|E| + B) \) using \( D \).
\end{itemize}


\end{problem}


\begin{solution}

    \

    \textbf{Data Structure: }

    \begin{itemize}
        \item \textbf{Setup:} We initialize an array of length equal to the maximum distance $B$. Initialize this array with null/empty linked list nodes.
        \item \textbf{Insert:} When we insert a new item into this priority queue, we find the bucket corresponding to its priority (between zero and the maximum B) and insert it into the linked list present there.
        \item \textbf{Extract-Min: } We traverse the array from beginning to end and return an item from the first non-null linked list(which is effectively a bucket).
        \item \textbf{Decrease-Key:} We find the vertex in the array(its index is equal to its priority. Then we link it to whatever corresponding bucket it belongs with(its new priority). 
    \end{itemize}

    \
    \textbf{Pseudocode: }

    \
    \begin{verbatim}
        def insert(int priority, int id):
            queue[priority].insert(id) # run a linked list insert on the linked list at the corresponding index for the priority of the item being added

        def extract_min(prev_min):
            for i in range(prev_min, queue.length):
                if queue[i] is not null:
                    return queue[i].pop()
        def decrease_key(int id, int old_priority, int new_priority): # we must get this information before deciding to update a key, so we know we have it when its time to change the key
            queue[old_priority].remove(id)
            queue[new_priority]].insert(id)

    \end{verbatim}

    \textbf{Runtime: }
        Inserting and decreasing keys are constant time operations because they just access the array and insert into a linked list. 
        Because extract-min knows what the previous minimum is, although each run is O(B), accross all runs of extract min for a given graph, it will only take up O(B) time
    \

    \textbf{Runtime 2: }
        As discussed above, extract-min will use O(B) time for the life of the algorithm.

        The insert and decrease key will be run a maximum of V and E times respectively. This makes V irrelevant.

        In total, we now have a runtime of O(|E| + B)
\end{solution}

\newpage


\begin{problem}
Describe and analyze a modification of the Bellman–Ford algorithm that achieves the following:
\begin{itemize}
    \item If a negative cycle is reachable from \( s \):
    \begin{itemize}
        \item Return the negative cycle.
        \item Set \( \text{dist}(v) = -\infty \) for any vertex \( v \) reachable through the cycle.
    \end{itemize}
    \item If no negative cycle is reachable:
    \begin{itemize}
        \item Return the shortest-path tree.
        \item Compute the correct shortest-path distances \( \text{dist}(v) \) from \( s \) to every vertex \( v \).
    \end{itemize}
    \item The algorithm should run in \( O(VE) \) time.
\end{itemize}


\end{problem}


\begin{solution}
      
\end{solution}

\newpage



\begin{problem}
    Let $G=(V, E)$ be a directed graph with weighted edges, edge weights can be positive, negative, or zero. Suppose vertices of $G$ are partitioned into $k$ disjoint subsets $V_1, V_2, \ldots, V_k$; that is, every vertex of $G$ belongs to exactly one subset $V_i$. For each $i$ and $j$, let $\delta(i, j)$ denote the minimum shortest-path distance between vertices in $V_i$ and vertices in $V_j$, that is

$$
\delta(i, j)=\min \left\{\operatorname{dist}(u, v) \mid u \in V_i \text { and } v \in V_j\right\}
$$


Describe an algorithm to compute $\delta(i, j)$ for all $i$ and $j$ that runs in $O(V E+k V \log V)$ time. $O(V^3)$ algorithm will get partial credit.
\end{problem}



\begin{solution}

%=================================================
    % 1. THE ALGORITHM %=====================================================





%=================================================
    % 2. PROOF OF CORRECTNESS %=====================================================





%=================================================
    % 3. RUN TIME ANALYSIS %=====================================================


\end{solution}

\newpage







\end{document}
